\section{Study Design}
\label{sec:setup}

\subsection{Pipeline, schemas and enforcement}
\label{sec:setup:pipeline}

Each run processes one question through a three-node sequential pipeline.
\textbf{identify} extracts entities with canonical identifiers and types
and designates a primary entity; \textbf{gather} produces facts keyed to
those identifiers with a coverage summary; \textbf{answer} returns the
answer, a numeric confidence, its categorical bucket and supporting
identifiers. Each node has its own Pydantic model. The system prompt
states the schema and asks for JSON only; the completion is validated
with \code{model\_validate\_json}; on failure the node is retried with the
validation error appended, up to three further times. An \emph{attempt} is
one generation against one node's schema, including retries, so attempts
per run vary with failure rate: from
\pn{c.claude-opus-4-8@medium@hotpotqa.attempts_per_run} for Opus~4.8 to
\pn{c.claude-haiku-4-5@medium@hotpotqa.attempts_per_run} for Haiku~4.5,
both on HotpotQA.

The \emph{medium} schema tier carries the constraints analysed
throughout: the identifier pattern \idpattern{} on entity and fact
identifiers, enumerations on entity and question type, character caps
on free text (\code{coverage\_notes} 250, \code{answer} 400, \code{fact}
400), list-length bounds and numeric ranges. Model-level validators add
four cross-field checks---primary-entity membership and entity-count
agreement in identify, fact-count agreement in gather, and
confidence-bucket coherence in answer---and, in answer, a per-item
identifier pattern check on \code{supporting\_entity\_ids}. The system
prompt carries each node's JSON Schema, so field-level constraints and
field descriptions are visible from the first attempt. The four
cross-field checks are stated in field descriptions; the
supporting-identifier check is stated nowhere in the prompt and becomes
visible only through retry feedback (\cref{tab:visibility}).
Two targeted variants of the baseline schema each change one constraint
family. \emph{medium\_v2} raises the caps to 1000, 1200 and 800 and is
run on Sonnet~4.6 only; it also omits the three node-level docstrings,
which the prompt carries as schema descriptions but which state no
constraint. \emph{medium\_v3\_loose\_regex} loosens the identifier
pattern to \idpatternloose{}, keeping the 31-character ceiling, rewords
the \code{canonical\_id} description to match, and is run on
gpt-5.6-sol and gpt-5.6-terra only. The design is therefore not fully
crossed; \cref{app:manifest} lists every cell.

\subsection{Models, data and reporting basis}
\label{sec:setup:basis}

Nine models from two vendors: \code{gpt-4o}, \code{gpt-4o-mini},
\code{gpt-5.6-luna}, \code{gpt-5.6-sol}, \code{gpt-5.6-terra},
\code{claude-haiku-4-5-20251001}, \code{claude-sonnet-4-6},
\code{claude-sonnet-5} and \code{claude-opus-4-8}, all called through the
vendors' own APIs between \pn{collection.date_min} and
\pn{collection.date_max}. Only the Haiku request names a dated snapshot;
the two gpt-4o models resolved to dated snapshots, and the remaining
endpoints returned the requested alias, which identifies the endpoint but
not an immutable model version (\cref{app:manifest}). HotpotQA
\citep{yang2018hotpotqa} questions are the first records of the
validation split (distractor configuration) and FinanceBench
\citep{islam2023financebench} questions the first records of the public
sample released as \code{PatronusAI/financebench}, which is a subset of the
full FinanceBench question set and has a single \code{train} split, with
evidence truncated to 2{,}000 characters; neither is shuffled,
so results describe these fixed question sets. Most cells contain about
one run per question; the Haiku cells repeat questions, and its
FinanceBench cell covers
\pn{c.claude-haiku-4-5@medium@financebench.questions} distinct questions
(\cref{app:manifest}). Claude Sonnet~4.6 and gpt-4o were run on HotpotQA
only. We request \code{temperature=0} where supported; the GPT-5.6 family,
Sonnet~5 and Opus~4.8 reject the parameter and run at provider defaults.
The Anthropic branch sets \code{max\_tokens} to 4{,}096 and the OpenAI
branch sets none. Resolved model versions and sampling details are in
\cref{app:manifest}.

Two conventions define the reporting basis. Cells with fewer than 30
baseline runs are excluded as pilots, a threshold fixed before analysis,
and the \pn{excluded_variant.runs} prompt-variant runs, which share the
Haiku HotpotQA configuration, are excluded everywhere except
\cref{sec:results:concealment}. The basis contains
\pn{reported.attempts} attempts and \pn{reported.failures} failures in
\pn{reported.cells} configurations, from a full corpus of
\pn{full.attempts} attempts in \pn{full.cells} cells. Classifier
validation (\cref{sec:setup:taxonomy}) uses the full corpus and says so.

\paragraph{Uncertainty.}
Attempts are not independent: retries follow failures, questions repeat
in some cells, and later nodes are reached only after earlier ones
succeed. Intervals are therefore from a question-clustered bootstrap
(\pn{meta.bootstrap_reps} replicates, percentile 95\%) that resamples
questions within each dataset and carries every run and attempt of a
selected question. They describe variability over the observed question
set and do not make the fixed prefixes representative of the benchmarks.
Because a failure adds an attempt under retry-until-valid, failures and
attempts co-vary within a question, so for high-failure models the
clustered interval on the failure rate can be narrower than a binomial
interval over attempts. Where every replicate gives the same value, as for
a share of 0 or 1, we report no interval.

\subsection{Taxonomy and classifier}
\label{sec:setup:taxonomy}

The classifier runs only on attempts that failed validation and assigns
one primary category in two stages. \emph{Stage~1} applies behavioural
predicates to the raw completion $y$, in order: \cat{markdown\_wrap} if
$y$ begins, after leading whitespace, with a triple-backtick fence;
\cat{prose\_preamble} if $y$ begins with neither a brace, a bracket nor a
fence, contains a brace-delimited span, and has at least 20 characters of
prose before it; and \cat{schema\_mirroring} if $y$ parses to an object
whose root keys include at least two JSON~Schema meta-keywords such as
\code{properties} and \code{required}. \emph{Stage~2}, reached only if no
predicate fires, maps Pydantic error types: pattern mismatches to
\cat{regex\_pattern}, length violations to \cat{length\_bound}, literal
errors to \cat{enum\_violation}, missing fields to \cat{missing\_field},
numeric bounds to \cat{range\_bound} and \code{json\_invalid} to
\cat{parse\_error}. Type errors map to \cat{range\_bound} for numeric
fields and to \cat{missing\_field} otherwise. Validator errors arrive as \code{value\_error}; the
supporting-identifier check routes to \cat{regex\_pattern} and cross-field
messages to \cat{cross\_field\_ref}. With several errors, the primary
category is the plurality, ties broken by first occurrence.

We give Stage-1 predicates precedence to preserve behaviours that are
invisible to Pydantic, which sees only their downstream symptoms. The
precedence has a consequence we exploit in
\cref{sec:results:concealment}: a fenced output that also violates a
constraint is labelled \cat{markdown\_wrap}, so that category measures
strict-parsing failure, not the absence of a further defect.
\Cref{tab:categories} gives the observed distribution;
\pn{reported.categories_observed} of the classifier's eleven categories
occur, and \cat{cross\_field\_ref}, \cat{missing\_field},
\cat{range\_bound} and the catch-all category do not.

\begin{table}[t]
\centering\small
\caption{Primary categories on the reporting basis.}
\label{tab:categories}
\begin{tabular}{@{}lrr@{}}
\toprule
Category & $n$ & Share \\
\midrule
\cat{markdown\_wrap} & 681 & 60.5\% \\
\cat{length\_bound} & 263 & 23.4\% \\
\cat{regex\_pattern} & 126 & 11.2\% \\
\cat{schema\_mirroring} & 38 & 3.4\% \\
\cat{prose\_preamble} & 9 & 0.8\% \\
\cat{enum\_violation} & 8 & 0.7\% \\
\cat{parse\_error} & 1 & 0.1\% \\
\midrule
Total & 1{,}126 &  \\
\bottomrule
\end{tabular}

\end{table}

\paragraph{Validation.}
We hand-labelled \pn{validation.n} failing attempts stratified by
predicted category, blind to the classifier's output, drawn from the full
corpus of \pn{full.failures} failures so that rare categories had
support. Labelling surfaced a detection-order defect in an earlier
classifier, in which Pydantic mapping preceded the behavioural
predicates; we corrected the order and reclassified the corpus. Against
the corrected classifier, blind agreement is \pn{validation.blind_agree}
of \pn{validation.n} (\pn{validation.blind_rate}, Wilson 95\% CI over labelled attempts
\pn{validation.blind_ci}). On adjudication all four disagreements were
labelling errors. Because these samples also surfaced the defect, this is
development-set agreement, not held-out validation, and a single labeller
who also wrote the classifier produced every label.
\Cref{app:classifier} gives per-category precision and recall and
over-broad corpus screens for the Stage-1 categories.